\chapter{Localization of the lowest quantiles}\label{sec:lz-root}

In this section, we show that the lowest quantiles of the free mass of every
minimizer of $F_\rho$ lie close to the origin. Let $\mu=s\delta_0+\nu$ be a
minimizing profile and $\eta$ a trial budget. By \cref{cor:lz-omega}, the
quantile $Q(v)$ of $\nu$ at level $v$ is at most
$C_Q(s+v+\omega(\rho,\eta))$, where
$\omega(\rho,\eta)=\min\{1,(\eta/\rho)^{1/6}\}$ and $C_Q$ depends only on
$\beta$. \Cref{prop:lz-quantile} is the same estimate in the energy
normalization. In \cref{prop:sm-small}, we use \cref{cor:lz-omega} to bound the
self-overlap variance on dyadic quantile shells for the interpolation
parameters near $\rho=0$.

We keep the standing assumptions and the notation of \cref{subsec:lz-inputs}.
From \cref{sec:lz}, we use the convexity and the first variation of $F_\rho$
(\cref{prop:loc-inputs}(a)), the identities \eqref{eq:loc-b}, the coercivity
bound of \cref{lem:loc-coercive}, and the midpoint inequality
\eqref{eq:loc-midpoint}. We also use the partial regularization of
\cref{lem:loc-smoothing}(a), the scalar root bounds \eqref{eq:lz-gsroot}, the
structure \eqref{eq:loc-scalar-measure} of $\mu_\beta$, and the trial budget
\eqref{eq:loc-trialbudget}. The isolation estimate \cref{prop:loc-isolation}
and the density bound \cref{lem:loc-lipschitz} are not used here.

The main tool is a local deficit certificate (\cref{lem:lz-certificate}). It
combines a variation that moves mass of a
profile toward the middle of an interval with the response inequality
\eqref{eq:loc-response}. If the variation violates its room or sign
condition, this failure gives a lower bound on $\eta/\rho$ as well, and the
certificate includes this bound as an alternative. In
\cref{subsec:lz-quantile}, we follow the strategy of the proof of
\cite[Lemma~7.6]{AL}. We regularize only the free mass below the selected
quantile. Then we compare the result with the scalar trial profile on an
interval where the latter has a positive density (\cref{lem:lz-rootdensity}).
The certificate gives the dependence on $\eta/\rho$ in \eqref{eq:lz-quantile}.

\section{A local deficit certificate}\label{subsec:lz-moves}

In the proof of \cref{prop:loc-isolation}, we used the uniform lower bound
\eqref{eq:loc-in-curv} on $\zeta$, which gives
$\Delta=\rho\int\zeta\varphi\ge\rho k_\beta^2\int\varphi$. For root
localization, we need $\Delta$ to be at least a constant multiple of the length
of the interval on which $\varphi$ lives, without the factor $\rho$. In the
next lemma, we achieve this by moving mass of the profile toward the middle of
the interval. If the variation does not have enough room, or if it would
destroy the profile gap, then one of the lower bounds in the conclusion already
holds.

\begin{lemma}[Local deficit certificate]\label{lem:lz-certificate}
There exists $c_{\rm cert}=c_{\rm cert}(\beta)>0$ with the following property.
Let $I\Subset(0,1)$ be an open interval of length $\ell$, let
$\gamma_0,\gamma_1$ be admissible profiles, and let $\eta>0$ be such that
$F_\rho(\gamma_i)\le\min F_\rho+\eta$ for $i=0,1$. Suppose that there exist
$k>0$, $d_0,d_1\ge0$, and $A>0$ such that, on $I$, the derivative measures and
the profile gap satisfy
\[
 \dd\gamma_0\le d_0\dd x,\qquad
 k\dd x\le\dd\gamma_1\le d_1\dd x,
 \qquad \gamma_1-\gamma_0\ge A .
\]
Then
\begin{equation}\label{eq:lz-certificate}
 \frac\eta\rho\ge c_{\rm cert}\min\left\{
 k\ell^3,\ A\ell^2,\ A^2\ell^4,
                  \frac{A^2\ell^2}{d_0+d_1+k}\right\}.
\end{equation}
\end{lemma}

\begin{proof}
The proof has four steps. In Step 1, we fix a bump function $\varphi$ on $I$
and an amplitude $\tau$. If $\tau$ is too large for the variation below, then
one of the first two alternatives in \eqref{eq:lz-certificate} holds. Otherwise, we perturb the segment from
$\gamma_0$ to $\gamma_1$ by the derivative of the bump, which moves mass of
the profile toward the middle of $I$ (Step 2). Since both endpoints are nearly
minimal and $F_\rho$ is convex, the response integrated against the bump is
at least of order $\ell$ along the perturbed segment (Step 3). Then
\cref{lem:loc-coercive} implies a lower bound on the curvature along the
perturbed segment. Comparing this curvature with the excesses at its
endpoints, we obtain the last two alternatives (Step 4). Throughout, we write $F_{\min}=\min F_\rho$ and
$I=(x_0,x_0+\ell)$.

\paragraph{Step 1: The bump and the room condition.}
Let $B(y)=16y^2(1-y)^2$ for $y\in[0,1]$, and let $\varphi(x)=B((x-x_0)/\ell)$
for $x\in I$ and $\varphi(x)=0$ for $x\notin I$. We have
\[
 B'(y)=32y(1-y)(1-2y),\qquad B''(y)=32-192y+192y^2,
\]
and $B$ and $B'$ vanish at $0$ and $1$. A direct computation shows that
$\int_0^1B=8/15$, $\max|B'|=16\sqrt3/9<4$, $\max|B''|=32$, and
$\int_0^1|B'|^2=512/105$. These properties of $B$ imply that $\varphi$ is
continuously differentiable on $\R$ with values in $[0,1]$, that $\varphi'$ is
Lipschitz, that $\varphi$ and $\varphi'$ vanish at the endpoints of $I$, and
that
\begin{equation}\label{eq:lz-bump}
 \int\varphi=\frac{8\ell}{15},\qquad \|\varphi'\|_\infty\le\frac4\ell,\qquad
 \|\varphi''\|_\infty\le\frac{32}{\ell^2},\qquad
 \int|\varphi'|^2=\frac{512}{105\,\ell},
\end{equation}
where $\varphi''$ is defined away from the endpoints of $I$. In particular,
the distributional derivative of $\varphi'$ is the function $\varphi''$, which
has no atoms at the endpoints of $I$, and $\int\varphi''=0$.

We define $\tau=30\eta/(\rho\ell)$; the factor $30$ is chosen so that
$8\eta/(\rho\tau)=4\ell/15$ is half of $\int\varphi$. Suppose first that the
room condition
\begin{equation}\label{eq:lz-room-condition}
 \tau\le\min\Bigl\{\frac{k\ell^2}{64},\frac{A\ell}8\Bigr\}
\end{equation}
fails. Then
\[
 \frac\eta\rho=\frac{\ell\tau}{30}
 >\min\Bigl\{\frac{k\ell^3}{1920},\frac{A\ell^2}{240}\Bigr\}
 \ge\frac1{1920}\min\{k\ell^3,A\ell^2\},
\]
and \eqref{eq:lz-certificate} holds whenever $c_{\rm cert}\le1/1920$. We
assume from now on that \eqref{eq:lz-room-condition} holds.

\paragraph{Step 2: The perturbed segment.}
For $0\le\theta\le1$, let
\[
 \gamma_\theta=(1-\theta)\gamma_0+\theta\gamma_1,\qquad
 \widetilde\gamma_\theta=\gamma_\theta-\theta\tau\varphi'.
\]
We claim that $\gamma_\theta-u\varphi'$ is admissible for $0\le\theta\le1$ and
$0\le u\le\theta\tau$. The function $\varphi'$ is continuous and vanishes
outside $I$, and $\overline I\subset(0,1)$. Then $\gamma_\theta-u\varphi'$ is
right-continuous and agrees with $\gamma_\theta$ near $0$ and near $1$. Its
derivative measure is $\dd\gamma_\theta-u\varphi''\dd x$. Off $I$, this is
$\dd\gamma_\theta$. On $I$, the hypotheses give $\dd\gamma_\theta\ge\theta
k\dd x$, and by \eqref{eq:lz-bump} and \eqref{eq:lz-room-condition}, we have
\[
 \dd(\gamma_\theta-u\varphi')
 \ge\Bigl(\theta k-\frac{32u}{\ell^2}\Bigr)\dd x
 \ge\Bigl(\theta k-\frac{32\theta\tau}{\ell^2}\Bigr)\dd x
 \ge\frac{\theta k}2\dd x .
\]
The signed measure $\varphi''\dd x$ has total mass zero and does not charge
$0$. We conclude that the derivative measure of $\gamma_\theta-u\varphi'$ is
nonnegative, has total mass $\beta$, and has the same atom at $0$ as
$\dd\gamma_\theta$; that is, $\gamma_\theta-u\varphi'$ is admissible. We have
$\widetilde\gamma_\theta=(1-\theta)\gamma_0+\theta(\gamma_1-\tau\varphi')$, and
$\gamma_1-\tau\varphi'$ is admissible (this is the claim with $\theta=1$ and $u=\tau$). It
follows that $(\widetilde\gamma_\theta)$ is a segment, with endpoints
$\widetilde\gamma_0=\gamma_0$ and $\widetilde\gamma_1=\gamma_1-\tau\varphi'$.

\paragraph{Step 3: The response integrated against the bump.}
Fix $\theta\in(0,1]$ and define $g_\theta(u)=F_\rho(\gamma_\theta-u\varphi')$
for $0\le u\le\theta\tau$. The profiles $\gamma_\theta-u\varphi'$,
$0\le u\le\theta\tau$, form a segment up to a linear change of parameter, so
$g_\theta$ is convex and continuous by \cref{prop:loc-inputs}. We compute its
one-sided derivatives. Let $0\le u\le\theta\tau$, and write
$\gamma'=\gamma_\theta-u\varphi'$, $G=G_{\gamma'}$, and $S=S_{\gamma'}$.
Increasing $u$ changes the free measure of $\gamma'$ in the direction
$-\varphi''\dd x$, and decreasing $u$ changes it in the direction
$+\varphi''\dd x$; in both cases, the new free measure is that of an
admissible profile on the segment. By \eqref{eq:loc-firstvariation}, the right
derivative of $g_\theta$ at $u<\theta\tau$ and the left derivative at $u>0$
are both equal to $-(\rho/2)\int G\varphi''\dd x$. We integrate by parts
twice. The function $\varphi'$ vanishes at the endpoints of $I$ and
$G'(x)=x-S(x)$, which gives $\int G\varphi''\dd x=-\int(x-S)\varphi'\dd x$. The
function $\varphi$ also vanishes at the endpoints of $I$, and $S$ is absolutely
continuous by \eqref{eq:loc-b}; integrating by parts once more, we find
\[
 -\int(x-S)\varphi'\dd x=\int\varphi\dd(x-S)=\frac{8\ell}{15}-\int\varphi\dd S .
\]
We obtain
\begin{equation}\label{eq:lz-concentrating-derivative}
 g_\theta'(u)=\frac\rho2
          \left(\int\varphi\dd S_{\gamma_\theta-u\varphi'}-\frac{8\ell}{15}\right),
\end{equation}
for the right derivative at $0\le u<\theta\tau$ and for the left derivative
at $0<u\le\theta\tau$.

By the convexity of $F_\rho$ along the segment $(\gamma_\theta)$, we have
$g_\theta(0)=F_\rho(\gamma_\theta)\le(1-\theta)F_\rho(\gamma_0)
+\theta F_\rho(\gamma_1)\le F_{\min}+\eta$, and $g_\theta(\theta\tau)\ge
F_{\min}$. For a convex function, the left derivative at the right endpoint
of an interval is at least the slope of the secant over the interval, which
gives
\[
 g_\theta'(\theta\tau)\ge\frac{g_\theta(\theta\tau)-g_\theta(0)}{\theta\tau}
 \ge-\frac\eta{\theta\tau}.
\]
Inserting \eqref{eq:lz-concentrating-derivative} at $u=\theta\tau$, where
$\gamma_\theta-\theta\tau\varphi'=\widetilde\gamma_\theta$, we find, for
$1/4\le\theta\le3/4$,
\[
 \Delta_\theta=\int\varphi\dd S_{\widetilde\gamma_\theta}
 \ge\frac{8\ell}{15}-\frac{2\eta}{\rho\theta\tau}
 \ge\frac{8\ell}{15}-\frac{8\eta}{\rho\tau}=\frac{4\ell}{15}.
\]
By \eqref{eq:loc-b}, we have $\dd S_{\widetilde\gamma_\theta}
=\rho\zeta_{\widetilde\gamma_\theta}\dd x$. In other words,
$\Delta_\theta=\rho\int\zeta_{\widetilde\gamma_\theta}\varphi\dd x$ is the
quantity $\Delta$ of \cref{lem:loc-coercive} for the segment
$(\widetilde\gamma_\theta)$ and the cutoff $\varphi$.

\paragraph{Step 4: Curvature and the endpoint excesses.}
By \eqref{eq:lz-bump} and \eqref{eq:lz-room-condition}, the direction of the
segment $(\widetilde\gamma_\theta)$ satisfies
\[
 \widetilde\gamma_1-\widetilde\gamma_0=\gamma_1-\gamma_0-\tau\varphi'
 \ge A-\frac{4\tau}\ell\ge\frac A2\qquad\text{on }I=\{\varphi>0\}.
\]
On $I$, the derivative measure of $\widetilde\gamma_\theta$ is at most
$((1-\theta)d_0+\theta d_1+32\theta\tau/\ell^2)\dd x\le\bar d\dd x$, where
\[
 \bar d=d_0+d_1+\frac{32\tau}{\ell^2}\le d_0+d_1+\frac k2 ;
\]
note that $\bar d\ge d_1\ge k>0$. Let
$\tilde f(\theta)=F_\rho(\widetilde\gamma_\theta)$. By \cref{lem:loc-coercive},
used for the segment $(\widetilde\gamma_\theta)$ with the cutoff $\varphi$,
$\sigma=1$, the level $A$, and $d=\bar d$, together with
$\Delta_\theta\ge4\ell/15$ and $Y_0=512/(105\ell)$ from \eqref{eq:lz-bump}, we
obtain, for $1/4\le\theta\le3/4$,
\[
 \tilde f''(\theta)\ge\frac{\rho A^2}{72}
 \min\left\{\Bigl(\frac{4\ell}{15}\Bigr)^2\frac{105\,\ell}{512\beta},
 \frac{4\ell}{15\bar d}\right\}
 =\frac{\rho A^2}{72}\min\left\{\frac{7\ell^3}{480\beta},
 \frac{4\ell}{15\bar d}\right\}.
\]
Because $\beta>1$, we have $4/15\ge7/(480\beta)$, which gives
\begin{equation}\label{eq:lz-certificate-curvature}
 \frac{\dd^2}{\dd\theta^2}F_\rho(\widetilde\gamma_\theta)
 \ge\frac7{34560\beta}\,\rho A^2\min\Bigl\{\ell^3,\frac\ell{\bar d}\Bigr\}
 \qquad(1/4\le\theta\le3/4).
\end{equation}

It remains to bound the excesses at the endpoints of the perturbed segment.
At $\theta=0$, we have $\tilde f(0)=F_\rho(\gamma_0)\le F_{\min}+\eta$. At
$\theta=1$, we have $\tilde f(1)=g_1(\tau)$, and the convexity of $g_1$ gives
$g_1(\tau)-g_1(0)\le\tau g_1'(\tau)$, with the left derivative. By
\eqref{eq:loc-b}, the function $S_{\widetilde\gamma_1}$ is nondecreasing with
values in $[0,1]$, and since $0\le\varphi\le1$, we have
$0\le\int\varphi\dd S_{\widetilde\gamma_1}\le1$. Then
\eqref{eq:lz-concentrating-derivative} gives $g_1'(\tau)\le\rho/2$. We
conclude that $\tilde f(1)\le F_\rho(\gamma_1)+\rho\tau/2\le F_{\min}+\eta+15\eta/\ell$.
Using $\ell\le1$, we obtain
\[
 \tilde f(0)+\tilde f(1)-2F_{\min}\le2\eta+\frac{15\eta}\ell
 \le\frac{17\eta}\ell .
\]
We now use the midpoint inequality \eqref{eq:loc-midpoint} for the segment
$(\widetilde\gamma_\theta)$ with $\Theta=[1/4,3/4]$, on which the weight has
integral $3/16$. Together with \eqref{eq:lz-certificate-curvature}, it yields
\[
 \frac3{16}\cdot\frac7{34560\beta}\,\rho A^2
 \min\Bigl\{\ell^3,\frac\ell{\bar d}\Bigr\}\le\frac{17\eta}\ell,
\]
that is,
\[
 \frac\eta\rho\ge\frac7{3133440\beta}A^2\min\Bigl\{\ell^4,\frac{\ell^2}{\bar d}\Bigr\}
 \ge\frac7{3133440\beta}\min\Bigl\{A^2\ell^4,\frac{A^2\ell^2}{d_0+d_1+k}\Bigr\},
\]
where we used $\bar d\le d_0+d_1+k$. Together with Step 1, this proves
\eqref{eq:lz-certificate} with $c_{\rm cert}=7/(3133440\beta)$, which is
smaller than $1/1920$ because $\beta>1$.
\end{proof}

\section{Root localization}\label{subsec:lz-quantile}

For an admissible profile $\gamma$, we define the quantile of its free measure
in the energy normalization by
\[
 Q_\gamma(v)=\inf\{q\in[0,1]:\nu_\gamma([0,q])\ge v\},
 \qquad 0<v<\beta-\lambda .
\]
If a quantile of a minimizer is large compared with the prescribed mass plus
its level, then the minimizer is small on an interval where the scalar trial
profile has a positive density. We control the density of the partially
regularized minimizer there, and then compare the two profiles on that
interval using \cref{lem:lz-certificate}. We follow the strategy of the proof
of \cite[Lemma~7.6]{AL}, which regularizes only the free mass below the
selected quantile and compares the result with the scalar trial profile. We
combine this strategy with \cref{lem:lz-certificate}, which gives the
dependence on $\eta/\rho$ in \eqref{eq:lz-quantile}. The positive density of
the trial profile near the origin comes from the next lemma.

\begin{lemma}[Trial density near the root]\label{lem:lz-rootdensity}
There exists $T_1=T_1(\beta)\in(0,T_0]$ such that
\begin{equation}\label{eq:lz-rootdensity}
 \dd\gamma_\beta\ge\frac{k_0}2\dd x\qquad\text{as measures on }[0,T_1].
\end{equation}
\end{lemma}

\begin{proof}
Recall from \eqref{eq:loc-scalar-measure} that $\mu_\beta$ has no atom in
$[0,q_\beta)$, and that its density $\varrho_\beta$ is continuous on
$[0,q_\beta]$. We have $\gamma_\beta(x)=\beta\int_0^x\varrho_\beta$ for
$0\le x<q_\beta$, and $\gamma_\beta(x)/x\to\beta\varrho_\beta(0)$ as
$x\downarrow0$. By \eqref{eq:lz-gsroot}, we have $\gamma_\beta(x)/x\ge k_0$ for
$0<x\le T_0$, which implies that $\beta\varrho_\beta(0)\ge k_0$. By the
continuity of $\varrho_\beta$ at $0$, there exists
$T_1\in(0,\min\{T_0,q_\beta\})$ such that $\beta\varrho_\beta\ge k_0/2$ on
$[0,T_1]$. The measure $\dd\gamma_\beta$ has no atom at $0$, because
$\gamma_\beta(0)=0$, and on $[0,T_1]$ it equals $\beta\varrho_\beta\dd x$; this gives
\eqref{eq:lz-rootdensity}. The scalar Parisi measure is unique
\cite[Corollary~1]{AuffingerChen2015Uniqueness}, and $q_\beta$ and
$\varrho_\beta$ depend only on $\beta$. The constants $k_0$ and $T_0$ also
depend only on $\beta$, and we may therefore choose $T_1$ depending only on
$\beta$.
\end{proof}

The width $T_1$ is not explicit, because the continuity of $\varrho_\beta$ at
$0$ is a qualitative statement.

\begin{proposition}[Root localization]\label{prop:lz-quantile}
There exists $C=C(\beta)<\infty$ with the following property. Let $N\ge1$,
$0<\rho\le1$, and $0<\lambda\le\lambda_0$, let $\gamma_*$ be a minimizer
of $F_\rho$, and let $\eta$ be a trial budget for $\gamma_*$. Then
\begin{equation}\label{eq:lz-quantile}
 Q_{\gamma_*}(v)\le C\left[\lambda+v+\min\{1,(\eta/\rho)^{1/6}\}\right]
 \qquad(0<v<\beta-\lambda).
\end{equation}
The constant depends on $\beta$ only through $\beta$, $k_0$, $C_0$, and
$T_1$.
\end{proposition}

\begin{proof}
The proof has five steps. After reducing to $0<\eta<\rho$, we replace the
quantile by its truncation $q=\min\{Q_{\gamma_*}(v),T_1\}$ and dispose of the
case where $q$ is small compared with $\lambda+v$ (Step 1). Otherwise, we
regularize the free mass of $\gamma_*$ below $q$, with a width adapted to the
moved mass (Step 2), compare the result with $\gamma^{\rm tr}$ on the middle
interval $(q/4,3q/4)$ (Step 3), and apply \cref{lem:lz-certificate} (Step 4).
The resulting inequality bounds $q$ either by a power of $\eta/\rho$ or in
terms of the density of the regularized profile, which we control with the
choice of the width and Young's inequality (Step 5).

\paragraph{Step 1: Reductions.}
Fix $v\in(0,\beta-\lambda)$, and let $\bar v=\lambda+v$. The constant $C$
constructed below is at least one. If $\eta\ge\rho$, the claim follows from
$Q_{\gamma_*}(v)\le1=\min\{1,(\eta/\rho)^{1/6}\}$. If $\eta=0$, every
$\eta'\in(0,\rho)$ is also a trial budget for $\gamma_*$; we apply the
argument below for positive budgets to $\eta'$ and let $\eta'\downarrow0$. We
may therefore assume that $0<\eta<\rho$, and we write
$\varepsilon=\eta/\rho\in(0,1)$.

We set $q=\min\{Q_{\gamma_*}(v),T_1\}$. Then $Q_{\gamma_*}(v)\le q/T_1$:
if $Q_{\gamma_*}(v)\le T_1$, then $q=Q_{\gamma_*}(v)$ and $T_1\le1$; otherwise
$q=T_1$ and $Q_{\gamma_*}(v)\le1=q/T_1$. It suffices to bound $q$.
If $q\le16\bar v/k_0$, then $Q_{\gamma_*}(v)\le16\bar v/(k_0T_1)$, which is of
the required form. Suppose instead that
\begin{equation}\label{eq:lz-q-large}
 q>16\bar v/k_0 .
\end{equation}

\paragraph{Step 2: Regularization below $q$.}
Let $m=\nu_{\gamma_*}([0,q))$. For every $x<q$, we have $x<Q_{\gamma_*}(v)$,
and the definition of the quantile gives $\nu_{\gamma_*}([0,x])<v$. Since the
interval $[0,q)$ is the increasing union of the intervals $[0,x]$, $x<q$,
continuity from below implies that $m\le v$. An atom of $\nu_{\gamma_*}$ at $q$ is not counted
in $m$: it is not moved by the regularization below, and it lies outside the
closure of the interval of Step 3.

We take
\[
 h=\min\{q/8,\sqrt{\varepsilon/m}\}\quad\text{if }m>0,
 \qquad h=q/8\quad\text{if }m=0 .
\]
Then $h\le q/8$ and $mh^2\le\varepsilon$, and $h\le1-q$ because $q\le
T_1\le T_0\le1/4$. Let $\bar\gamma$ be the partial regularization of
\cref{lem:loc-smoothing}(a) with $T=q$. By that lemma, we have
$F_\rho(\bar\gamma)-F_\rho(\gamma_*)\le\rho mh^2/12\le\eta/12$, and the
restriction of $\dd\bar\gamma$ to $(0,q)$ has density at most $m/h$. For
$0\le x<q$, the free mass of $\bar\gamma$ in $[0,x]$ comes from the moved
mass. Using \eqref{eq:lz-q-large} in the last step, we find that
\[
 \bar\gamma(x)\le\lambda+m\le\bar v<\frac{k_0q}{16}
 \qquad(0\le x<q).
\]

\paragraph{Step 3: The middle interval.}
Let $I=(q/4,3q/4)$. For $x\in I$, \eqref{eq:lz-gsroot} and
\eqref{eq:lz-q-large} give $\gamma_\beta(x)\ge k_0x>k_0q/4>4\bar v\ge\lambda$,
which implies that $\gamma^{\rm tr}=\gamma_\beta$ on $I$. Using
$I\subset(0,T_1]$, \cref{lem:lz-rootdensity}, and the density bound for
$\dd\gamma^{\rm tr}$ on $(0,T_0]$ stated after \eqref{eq:lz-gsroot}, we find
that
\[
 \frac{k_0}2\dd x\le\dd\gamma^{\rm tr}\le C_0\dd x\quad\text{on }I .
\]
By Step 2, we also have $\dd\bar\gamma\le(m/h)\dd x$ on $I$ and
\[
 \gamma^{\rm tr}-\bar\gamma>\frac{k_0q}4-\frac{k_0q}{16}\ge\frac{k_0q}8
 \quad\text{on }I .
\]

\paragraph{Step 4: The certificate.}
We apply \cref{lem:lz-certificate} with the interval
$I=(q/4,3q/4)\Subset(0,1)$, the endpoints $\gamma_0=\bar\gamma$ and
$\gamma_1=\gamma^{\rm tr}$, and the parameters
\[
 \ell=\frac q2,\qquad A=\frac{k_0q}8,\qquad k=\frac{k_0}2,
 \qquad d_0=\frac mh,\qquad d_1=C_0 .
\]
Both endpoints have excess at most $\eta$ over $\min F_\rho=F_\rho(\gamma_*)$:
for $\bar\gamma$ by Step 2, and for $\gamma^{\rm tr}$ by
\eqref{eq:loc-trialbudget}. The four entries of the minimum in
\eqref{eq:lz-certificate} are
\[
 \begin{gathered}
 k\ell^3=\frac{k_0q^3}{16},\qquad A\ell^2=\frac{k_0q^3}{32},\qquad
 A^2\ell^4=\frac{k_0^2q^6}{1024},\\
 \frac{A^2\ell^2}{d_0+d_1+k}=\frac{k_0^2q^4}{256(m/h+C_0+k_0/2)}.
 \end{gathered}
\]
Since $q\le1$, the first two entries are at least $k_0q^6/32$. With
$c'=c_{\rm cert}\min\{k_0/32,k_0^2/1024\}$, we find
\begin{equation}\label{eq:lz-balance}
 \varepsilon\ge c'\min\Bigl\{q^6,\frac{q^4}{m/h+C_0+k_0/2}\Bigr\}.
\end{equation}

\paragraph{Step 5: Conclusion.}
If the minimum in \eqref{eq:lz-balance} is $q^6$, then
$q\le(\varepsilon/c')^{1/6}$. Otherwise, we have
$q^4\le\varepsilon(m/h+C_0+k_0/2)/c'$. By the choice of $h$, we have
\[
 \frac mh=\max\Bigl\{\frac{8m}q,\varepsilon^{-1/2}m^{3/2}\Bigr\}
 \le\frac{8m}q+\varepsilon^{-1/2}m^{3/2}
\]
if $m>0$, and $m/h=0$ if $m=0$. Using $m\le v\le\bar v$
and \eqref{eq:lz-q-large}, we have $8m/q<k_0/2$. When the minimum is the second entry, we therefore
obtain
\[
 q^4\le\frac{(C_0+k_0)\varepsilon+\varepsilon^{1/2}m^{3/2}}{c'},
 \qquad
 q\le(c')^{-1/4}\bigl((C_0+k_0)^{1/4}\varepsilon^{1/4}
                           +\varepsilon^{1/8}m^{3/8}\bigr),
\]
where we used $(x+y)^{1/4}\le x^{1/4}+y^{1/4}$ for $x,y\ge0$. By Young's
inequality with exponents $8/5$ and $8/3$, we have
\[
 \varepsilon^{1/8}m^{3/8}
 \le\tfrac58\varepsilon^{1/5}+\tfrac38m
 \le\varepsilon^{1/6}+m,
\]
where in the last inequality we used that $0<\varepsilon<1$. Together with
$\varepsilon^{1/4}\le\varepsilon^{1/6}$ and $m\le\bar v$, this shows that
$q\le(c')^{-1/4}((C_0+k_0)^{1/4}+1)(\bar v+\varepsilon^{1/6})$. Combining the
three cases (Step 1, and the two cases of \eqref{eq:lz-balance}), we conclude
that
\[
 q\le C_1(\bar v+\varepsilon^{1/6}),\qquad
 C_1=\max\bigl\{16/k_0,(c')^{-1/6},(c')^{-1/4}((C_0+k_0)^{1/4}+1)\bigr\}.
\]
Recalling that $Q_{\gamma_*}(v)\le q/T_1$ and that
$\varepsilon^{1/6}=\min\{1,(\eta/\rho)^{1/6}\}$ for $\eta<\rho$, we obtain
\eqref{eq:lz-quantile} with $C=\max\{1,C_1/T_1\}$. This constant depends only
on $\beta$ (through $c_{\rm cert}$), $k_0$, $C_0$, and $T_1$.
\end{proof}

We comment on the exponent $1/6$ in \eqref{eq:lz-quantile}. In the setting of
the proof above, the variation in the proof of \cref{lem:lz-certificate}
moves mass toward the middle of $I$ at a cost of at most $\rho\tau/2$, which
is of order $\eta/q$, and the excesses at the endpoints of the perturbed
segment are at most of order $\eta/q$. The response integrated against the
bump is at least $4\ell/15=2q/15$. By \eqref{eq:lz-certificate-curvature}, the
second derivative along the perturbed segment is bounded below by a multiple
of $\rho\min\{q^5,q^3/\bar d\}$, where $\bar d$ is the density bound in the
proof of \cref{lem:lz-certificate}. Comparing this lower bound with the excesses at the endpoints, we obtain
$\eta/\rho\gtrsim_\beta\min\{q^6,q^4/\bar d\}$, which leads to
\eqref{eq:lz-balance}. Its first alternative gives the exponent $1/6$. In the
other alternative, the choice of the width $h$ controls the density $m/h$
without losing the linear dependence of the bound on $v$.

\begin{corollary}[Quantile bound in probability normalization]
\label{cor:lz-omega}
Let $N\ge1$, $0<\rho\le1$, and $0<\lambda\le\lambda_0$, and let
$s=\lambda/\beta$. Let $\mu=s\delta_0+\nu$ be a minimizer of
$\calF_{1-\rho,s}$ in \eqref{eq:fnd-functional}, so that
$\gamma_*=\beta\mu([0,\cdot])$ minimizes $F_\rho$, and let $\eta\ge0$ be a
trial budget for $\gamma_*$. Let
\[
 Q(v)=\inf\{q\in[0,1]:\nu([0,q])\ge v\}\qquad(0<v<1-s)
\]
be the quantile function \eqref{eq:fnd-quantile} of $\nu$, and define
\begin{equation}\label{eq:lz-omega}
 \omega(\rho,\eta)=\min\{1,(\eta/\rho)^{1/6}\}.
\end{equation}
Then, with $C_Q=\max\{1,\beta C\}$ and $C$ the constant of
\cref{prop:lz-quantile},
\begin{equation}\label{eq:lz-Q}
 Q(v)\le C_Q\bigl(s+v+\omega(\rho,\eta)\bigr)
 \qquad(0<v<1-s).
\end{equation}
\end{corollary}

\begin{proof}
The free measure of $\gamma_*$ is $\nu_{\gamma_*}=\beta\nu$, and
$\nu_{\gamma_*}([0,q])\ge\beta v$ if and only if $\nu([0,q])\ge v$. Hence
$Q(v)=Q_{\gamma_*}(\beta v)$, and $0<\beta v<\beta-\lambda$ for
$0<v<1-s$. By
\cref{prop:lz-quantile} and $\lambda=\beta s$, we have
\[
 Q(v)\le C\bigl[\beta(s+v)+\omega(\rho,\eta)\bigr]\le\beta C\bigl(s+v+\omega(\rho,\eta)\bigr),
\]
where we used $\beta>1$. We obtain \eqref{eq:lz-Q}, because $\beta C\le C_Q$.
\end{proof}
